% ============================================================
% Template: 如何书写一个 Model
% ============================================================
% 复制此文件并重命名，例如：harmonic-oscillator.tex
% 然后在 model-drill.tex 中对应 Part 下添加：
%   \input{models/harmonic-oscillator}
% ============================================================

\section{模型名称（中文）}
\label{sec:model-name}

\begin{tipbox}[关键词标签]
关键词1 · 关键词2 · 关键词3 · 关键词4
\end{tipbox}

% --------------------------------------------------------
% 1. 模型定义框（必须精准）
% --------------------------------------------------------
\begin{modelbox}[模型：XXX]
\textbf{物理情境}：一句话描述该模型出现的物理场景。

\textbf{哈密顿量}：
\[
\hat{H} = \dots
\]

\textbf{关键参数}：
\begin{itemize}[nosep]
    \item $a$：参数1的物理意义与量纲
    \item $b$：参数2的物理意义与量纲
\end{itemize}

\textbf{边界条件 / 约束}：
\begin{itemize}[nosep]
    \item 条件1
    \item 条件2
\end{itemize}
\end{modelbox}

% --------------------------------------------------------
% 2. 关键公式框（必须能独立推导）
% --------------------------------------------------------
\begin{formulabox}[核心公式]
\begin{enumerate}[label=\textbf{(\arabic*)}]
    \item \textbf{公式名称}：
    \[
    E_n = \dots
    \]
    \texteq{适用条件：}

    \item \textbf{公式名称}：
    \[
    \psi_n(x) = \dots
    \]
    \texteq{归一化条件：}
\end{enumerate}
\end{formulabox}

% --------------------------------------------------------
% 3. 训练点框（考法变体）
% --------------------------------------------------------
\begin{drillbox}[训练点与考法变体]
\trainingpoint{1} \textbf{基础考法}：直接求解该模型的能谱或波函数。

\trainingpoint{2} \textbf{参数扰动}：在模型上添加微扰 $V' = \dots$，求一级修正。

\trainingpoint{3} \textbf{边界条件变化}：改变边界条件（如从无限深势阱变为有限深），求能级变化趋势。

\trainingpoint{4} \textbf{模型组合}：与另一个模型组合（如 $A+B$），求耦合系统的行为。

\trainingpoint{5} \textbf{逆向问题}：已知测量结果，反推模型参数。
\end{drillbox}

% --------------------------------------------------------
% 4. 解题技巧框（考场快速识别）
% --------------------------------------------------------
\begin{tipbox}[快速识别与解题技巧]
\begin{itemize}
    \item \examnote{识别标志}：看到题干中出现 ``关键词''，立刻联想到本模型。
    \item \examnote{无量纲化}：引入特征长度 $x_0 = \sqrt{\hbar/m\omega}$，可将方程化为标准形式。
    \item \examnote{对称性利用}：若势能具有反演对称性，波函数必有确定宇称，可只解半区间。
    \item \examnote{量纲检查}：最终结果的量纲必须是 $\dots$
\end{itemize}
\end{tipbox}

% --------------------------------------------------------
% 5. 易错点框（肌肉记忆避免）
% --------------------------------------------------------
\begin{errorbox}{常见失误与陷阱}
\begin{itemize}
    \item \textbf{边界条件遗漏}：在 $x=0$ 处忘记连接条件 $\psi(0^+)=\psi(0^-)$。
    \item \textbf{符号陷阱}：$\delta$ 势的系数正负决定了是否存在束缚态。
    \item \textbf{归一化因子}：球坐标系中的体积元 $r^2\dd r$ 容易遗漏。
    \item \textbf{简并处理}：能级简并时，微扰论必须在对角化简并子空间后进行。
\end{itemize}
\end{errorbox}

% --------------------------------------------------------
% 6. 物理图像与关联模型
% --------------------------------------------------------
\begin{insightbox}[物理图像与模型关联]
\textbf{物理图像}：

用一句话或一个比喻描述该模型的核心物理图像。

\textbf{与相似模型的对比}：
\begin{center}
\begin{tabular}{lll}
\toprule
\textbf{对比维度} & \textbf{本模型} & \textbf{相似模型 A} \\
\midrule
势能形式 & $V(x)\sim x^2$ & $V(x)\sim |x|$ \\
能谱特征 & 等间距 & 非等间距 \\
适用场景 & 小振动 & 强耦合 \\
\bottomrule
\end{tabular}
\end{center}

\textbf{推广方向}：
\begin{itemize}[nosep]
    \item 维度推广：$1D \to 2D \to 3D$
    \item 相互作用推广：单体 $\to$ 多体
    \item 相对论推广：非相对论 $\to$ 相对论
\end{itemize}
\end{insightbox}

% --------------------------------------------------------
% 7. 推导附录（可选，供自我检验）
% --------------------------------------------------------
\begin{derivationbox}[推导细节（自我检验用）]
若考场上需要独立推导某公式，步骤如下：
\begin{enumerate}
    \item 步骤1：\dots
    \item 步骤2：\dots
    \item 步骤3：\dots
\end{enumerate}
\end{derivationbox}

\vspace{1em}
